\section{Varför en analysator}
\label{sec:whyanalyzer}

Fram till nu har varje observation gjorts av en av de två halvorna. Din drivrutin säger att den
skickade en frame; hårdvarugruppens testbänk säger att kontrollern sände en. Båda kan ha fel på
samma sätt, eftersom båda är skrivna mot samma dokument av två personer som kan ha läst det lika
fel.

En \textbf{bussanalysator} är en tredje part som inte tror på någon av er. Den är en riktig
CAN-kontroller på bussen, och den avkodar det som faktiskt ligger på ledningarna.

Det är därför labben ligger sist och inte först: det är först nu det finns trafik att titta på, och
det är först nu svaret betyder något.

\begin{aside}
\note Analysatorn kvitterar dessutom dina frames. Det är värt att veta: en ensam nod på bussen får
aldrig någon kvittens, och \emph{en fullständig} CAN-kontroller hade då skickat om i all evighet.
Konstruktionen i den här kursen läser aldrig tillbaka ACK-luckan, så den märker ingenting --- men
skillnaden blir synlig i tracen så fort analysatorn kopplas in.
\end{aside}

\section{Att läsa en trace}
\label{sec:trace}

Ställ in analysatorn på samma bithastighet som konstruktionen, och kontrollera att termineringen
fortfarande är 120~\textohm{} i vardera änden --- inte tre motstånd för att det nu är tre noder.

Skicka en frame med känt ID och känd nyttolast från din nod, och leta upp den:

\begin{booktable}{@{}lL@{}}
\thead{Kolumn} & \thead{Vad du kontrollerar} \\ \midrule
Identifierare & Stämmer den med vad du skrev till \code{TX_ID}? \\
DLC & Stämmer den med \code{TX_DLC}? \\
Data & Stämmer bytesen, \emph{och i rätt ordning}? Det är här en förväxling av
  \code{TX_DATA_LO}/\code{HI} äntligen blir synlig. \\
Tidsstämpel & Hur ofta kommer de? Stämmer det med hur ofta din kod anropar \code{send()}? \\
\end{booktable}

Sedan åt andra hållet: skicka en frame \emph{från} analysatorn till din nod. Returnerar
\code{receive()} \code{true}? Stämmer innehållet? Och --- det som är lätt att glömma --- gick
\code{STATUS} bit 1 ned när du skrev \code{RX_ACK}?

\section{Att leta efter förenklingarna}
\label{sec:simplifications}

Konstruktionen är förenklad mot riktig CAN på flera punkter. Några av dem går att se utifrån, och
det är den intressantaste delen av labben.

\begin{booktable}{@{}lL@{}}
\thead{Förenkling} & \thead{Vad du kan leta efter} \\ \midrule
Ingen ACK-återläsning & Koppla bort analysatorn och sänd med bara din nod på bussen. Vad säger
  felflaggan? Vad \emph{skulle} en riktig nod ha gjort? \\
Inga felramar & Vad händer om analysatorn skickar en? \\
Ingen automatisk omsändning & En frame som gick fel skickas inte om av hårdvaran. Syns det? \\
Ingen mottagningskö & Låt analysatorn skicka två frames tätt. Kommer båda fram till din
  applikation? \\
\end{booktable}

\textbf{Att hitta en avvikelse är inte att hitta ett fel.} Uppgiften är att kunna säga \emph{vilken
dokumenterad förenkling} som förklarar vad du ser --- och vilka avvikelser som \emph{inte} har
någon sådan förklaring. Det är de senare som är intressanta, för båda klasserna.

\section{Bussbelastning}
\label{sec:busload}

Analysatorn visar bussbelastningen. Jämför med tidsbudgeten i \kapref{ch:spi}:

\begin{itemize}
  \item Vad är belastningen med din nod som enda sändare?
  \item En \code{send()} tar sex SPI-transaktioner, alltså omkring 270~µs. En maximal CAN-frame är
    omkring 130~µs. \textbf{Vad begränsar dig egentligen --- bussen eller länken?}
\end{itemize}

Svaret på den sista frågan är hela poängen med att mäta i stället för att anta, och det är det
enda stället i kursen där du kan se det med egna ögon.

\section{En frame, hela vägen}
\label{sec:endtoend}

Avsluta med att spåra en enda frame genom hela stacken, och namnge varje lager och vem som byggde
det:

\begin{console}
app::EchoNode            -> du
driver::can::Interface   -> du
driver::can::Spi         -> du
driver::can::ByteTransport -> du
AvrSpiTransport          -> du
===== fyra ledningar =====
spi_slave                -> utdelad
spi_reg_bridge           -> hårdvarugruppen
register_bank            -> hårdvarugruppen
can_controller           -> hårdvarugruppen
===== transceiver =====
CAN-bussen               -> fysik
\end{console}

Tio lager, två klasser, ett kontrakt --- och en frame som kom fram.
